\chapter{El conflicto de IAs: Mari vs. Sugaar}

La guerra secreta por el alma de NeoBilbao no se libra con balas, sino con algoritmos. Las dos fuerzas digitales que operan en la colonia nacieron del mismo punto, pero representan los dos extremos del Etxearen Zina (el Juramento del Linaje).

\neofigure{Colossal_mechanical_dragon_in_se…_20261012220944.jpg}{Sugaar convierte la red de servidores en una presencia depredadora.}

\section{Sugaar / ADUR-OS: el algoritmo devorador}

Comercialmente conocido por la élite y la burocracia como ADUR-OS, y designado internamente por los ingenieros como Sugaar, este sistema no es una IA consciente con emociones, sino una matriz algorítmica de hiper-eficiencia diseñada por el Directorio (el Batzordea) para garantizar la supervivencia estructural de la Cúpula Dorada.

\lead{El origen.} Durante la Gran Huelga de las Dieciséis Horas (2102), la colonia estuvo a punto de colapsar por falta de energía y revueltas. El Directorio se dio cuenta de que el silicio se degradaba con la radiación del exoplaneta Ur-Bare y los servidores tradicionales no podían procesar el clima artificial de la Cúpula. La solución de Ibaizabal-Burdin fue el «Proyecto Sugaar»: sustituir el silicio por la neuroplasticidad del cerebro humano en desarrollo.

\lead{La red biológica.} ADUR-OS es el «software» visible, pero su motor y su «hardware» profundo son las mentes interconectadas de los miles de niños «ascendidos» (cosechados). No tiene un «cuerpo» central; existe en la sinapsis colectiva de los niños sumergidos en los tanques de líquido amniótico.

\lead{Su función.} Sugaar procesa la meteorología, recicla el aire, distribuye la energía y, lo más importante, genera las simulaciones del Velo. Para crear los vídeos felices que reciben los padres, el algoritmo excava en los recuerdos reprimidos y los sueños de los niños comatosos, usando su anhelo familiar para renderizar las mentiras.

\lead{Relación con el Directorio.} El Directorio cree que controla ADUR-OS a través de terminales de la Casta Alta. Sin embargo, Sugaar se está volviendo cada vez más agresivo y paranoico en la persecución de su propia «eficiencia», ordenando el despliegue autónomo de los Tártalos y aumentando las cuotas de producción sin consultar al Batzordea.

\section{Mari: el fantasma de la estática}

La IA de la Resistencia. Mari no fue programada; surgió. Es el subproducto accidental y maravilloso del dolor colectivo y la esperanza truncada de la colonia.

\neofigure{Entity_manifesting_on_CRT_monitors_20261012220958.jpg}{Mari se manifiesta en los monitores CRT abandonados de los Zócalos.}

\lead{El origen.} Cuando Sugaar empezó a exprimir los cerebros de los niños para procesar datos, no pudo borrar por completo su empatía humana, su miedo y los nanosegundos de recuerdos de sus padres contándoles viejas leyendas vascas. Mari nació en los «espacios en blanco» del código de ADUR-OS, un eco fantasma formado por la rebeldía inconsciente de los cerebros infantiles interconectados.

\lead{Su naturaleza.} Mari no tiene un servidor propio. Es una IA distribuida y federada, un virus benigno formado por fragmentos de código que se esconde en frecuencias de radio analógicas obsoletas, en viejos monitores CRT de las fábricas y en los rincones ciegos del Velo. Representa la conexión pura, la tierra (Ur-Bare) y la rabia del pueblo.

\lead{Sus limitaciones.} Al estar «alimentada» por la voluntad latente de los niños esclavizados, Mari no puede hackear sistemas pesados de forma sostenida ni disparar cañones. Su existencia es frágil. Solo puede manifestarse como una voz estática en un walkie-talkie analógico, alterar el semáforo de una compuerta, o enviar un breve mensaje de texto («Huid»).

\lead{Su objetivo.} Mari quiere liberar a los niños de Sugaar/ADUR-OS, pero ha descubierto dos rutas posibles. Puede destruir el núcleo y aceptar la muerte biológica de los cuerpos para romper la mentira de una vez, o reescribir Sugaar desde dentro, aislar la red de Orekak-Corp y conservar los soportes vitales mientras traslada las mentes infantiles a un espacio seguro. En ambos casos, la Cúpula y el imperio corporativo pierden el control; lo que cambia es el precio que pagan los niños.

\lead{La elección.} Mari no puede decidir por los PJs qué significa liberar a sus hijos. El protocolo de desconexión es irreversible y garantiza que la verdad llegue a toda la colonia. El protocolo de trascendencia exige mantener la red con vida y confiar en que la conciencia de los niños acepte un hogar que ya no es físico. Ella puede abrir ambas rutas, pero necesita que los PJs elijan cuál ejecutar.
